% Source: https://artofproblemsolving.com/community/c3h2789083_2011_manhattan_mathematical_olympiad_grades_78_problems_14
% Stable identity: aops-topic-2789083-q02
Let each tomato have radius $r$; the tomatoes occupy the same volume in both cans, so compare the cans.
\begin{asy}
usepackage("amsmath");
size(5cm);
defaultpen(fontsize(10pt));
draw((0,0)--(0,7)--(1,7)--(1,0)--cycle);for(int i=0;i<7;++i) draw(circle((0.5,0.5+i),0.5));label("tall can",(0.5,7.1),N);
pair c=(6.5,1.5);draw(circle(c,1.5));draw(circle(c,0.5));for(int k=0;k<6;++k) draw(circle(c+dir(60k),0.5));label("flat can, from above",(6.5,3.1),N);
\end{asy}

The tall can has radius $r$ and height $14r$: volume $14\pi r^3$. The flat can must hold a hexagonal ring of six tomatoes around one, so its radius is $3r$ and its height $2r$: volume $18\pi r^3$. The $\boxed{\text{wide flat can}}$ needs more brine (by $4\pi r^3$).
